%! Piecewise \& Step Functions — Unit 1, Lesson 1.4 — Group Activity
\documentclass[10pt]{article}
\usepackage{algebra2-article}
\usepackage{algebra2-boxes}
\newcommand{\TallMath}[1]{$\displaystyle #1\rule[-1.4em]{0pt}{3.2em}$}
% Coordinate grid: {xmin}{xmax}{ymin}{ymax}
\newcommand{\lgrid}[4]{%
  \draw[step=1, linegray!40, very thin] (#1,#3) grid (#2,#4);
  \draw[->, semithick, charcoal] (#1,0)--(#2,0) node[below right, font=\scriptsize]{$x$};
  \draw[->, semithick, charcoal] (0,#3)--(0,#4) node[left, font=\scriptsize]{$y$};}
% Endpoint dots
\newcommand{\closeddot}[2]{\fill[forest] (#1,#2) circle (3.5pt);}
\newcommand{\opendot}[2]{\draw[very thick, forest, fill=white] (#1,#2) circle (3.5pt);}

\begin{document}

\pageheader{Unit 1, Lesson 1.4}{Group Activity}

\vspace{0.08in}

\begin{headlinebox}{forestbg}
\small\textbf{Pieces, Graphs \& Staircases.} A piecewise function follows different rules on different
parts of its domain. Work with your partner. To \emph{evaluate}, first find the piece whose condition
your input satisfies, then apply that rule --- and watch the boundary ($<$ vs.\ $\le$). On a graph, a
\textbf{closed} dot $\bullet$ includes its endpoint and an \textbf{open} dot $\circ$ excludes it.
\end{headlinebox}

\vspace{0.06in}

\begin{tcolorbox}[colback=white, colframe=black!40, title=\textbf{Tier R --- Remediate},
    fonttitle=\bfseries, arc=1.5mm, left=3mm, right=3mm, top=2mm, bottom=2mm, breakable]
\begin{enumerate}[leftmargin=1.7em, itemsep=9pt]
  \item \textbf{Evaluate} $f(x)=\begin{cases} x+5, & x<0 \\ 2x, & x\ge 0 \end{cases}$ \quad at each input.
        \\[3pt]
        $f(-3)=\blank{1.3cm}$ \qquad $f(0)=\blank{1.3cm}$ \qquad $f(4)=\blank{1.3cm}$
  \item \textbf{Read the graph} of $g(x)=\begin{cases} x+1, & x<1 \\ 3, & x\ge 1 \end{cases}$.
        \begin{center}
        \begin{tikzpicture}[scale=0.5]
          \lgrid{-1}{4}{-1}{4}
          \draw[very thick, forest] (-1,0)--(1,2);
          \opendot{1}{2}
          \draw[very thick, forest] (1,3)--(4,3);
          \closeddot{1}{3}
        \end{tikzpicture}
        \end{center}
        $g(0)=\blank{1.0cm}$; \quad $g(1)=\blank{1.0cm}$; \quad $g(3)=\blank{1.0cm}$; \quad
        the \emph{closed} dot is at $(\blank{0.7cm},\,\blank{0.7cm})$.
  \item \textbf{Open or closed?} When you graph a piece\ldots
        \begin{itemize}[leftmargin=1.5em, itemsep=2pt]
          \item defined for $x>3$, the endpoint at $x=3$ is drawn \blank{2.2cm} (open / closed).
          \item defined for $x\le 3$, the endpoint at $x=3$ is drawn \blank{2.2cm} (open / closed).
        \end{itemize}
\end{enumerate}
\end{tcolorbox}

\begin{tcolorbox}[colback=white, colframe=black!40, title=\textbf{Tier A --- Approaching Proficiency},
    fonttitle=\bfseries, arc=1.5mm, left=3mm, right=3mm, top=2mm, bottom=2mm, breakable]
\begin{enumerate}[leftmargin=1.7em, itemsep=10pt]
  \item \textbf{Evaluate across the boundary.} For $h(x)=\begin{cases} -x+1, & x\le 2 \\ 2x-6, & x>2 \end{cases}$:
        \\[3pt]
        $h(0)=\blank{1.2cm}$ \quad $h(2)=\blank{1.2cm}$ \quad $h(5)=\blank{1.2cm}$. \quad
        Is $h$ continuous at $x=2$? \blank{1.4cm} \; (top gives \blank{0.8cm}, bottom gives \blank{0.8cm}).
  \item \textbf{Read the graph} below.
        \begin{center}
        \begin{tikzpicture}[scale=0.5]
          \lgrid{-2}{5}{-3}{3}
          \draw[very thick, forest] (-2,-2)--(2,2);
          \opendot{2}{2}
          \draw[very thick, forest] (2,-2)--(5,-2);
          \closeddot{2}{-2}
        \end{tikzpicture}
        \end{center}
        Domain \blank{2.2cm}; \; $f(1)=\blank{0.9cm}$; \; $f(2)=\blank{0.9cm}$; \; increasing on
        \blank{1.8cm}; \; constant on \blank{1.8cm}. \; Continuous at $x=2$? \blank{1.4cm}.
  \item \textbf{Write it as pieces.} Rewrite $|x-3|$ as a two-piece function.
        \[ |x-3|=\begin{cases} \rule{2.2cm}{0.4pt}, & x\ge 3 \\[3pt] \rule{2.2cm}{0.4pt}, & x<3 \end{cases} \]
  \item \textbf{Greatest integer.} $\lfloor 2.9\rfloor=\blank{1.0cm}$ \quad
        $\lfloor -1.2\rfloor=\blank{1.0cm}$ \quad $\lfloor 4\rfloor=\blank{1.0cm}$.
\end{enumerate}
\end{tcolorbox}

\begin{tcolorbox}[colback=white, colframe=black!40, title=\textbf{Tier E --- Extension},
    fonttitle=\bfseries, arc=1.5mm, left=3mm, right=3mm, top=2mm, bottom=2mm, breakable]
\begin{enumerate}[leftmargin=1.7em, itemsep=10pt]
  \item \textbf{Model overtime pay.} A worker earns \$15/hr for up to $40$ hours, then \$22.50/hr
        (time-and-a-half) for every hour beyond $40$. Let $P(h)$ be the weekly pay for $h$ hours.
        \begin{enumerate}[label=(\alph*), leftmargin=1.8em, itemsep=5pt]
          \item Complete the rule:
                $P(h)=\begin{cases} \rule{2.0cm}{0.4pt}, & 0\le h\le 40 \\[3pt] 600+\rule{2.6cm}{0.4pt}, & h>40 \end{cases}$
          \item Find $P(40)$ and $P(45)$, and say what $P(45)$ means.
                \begin{work}
                  P(40) &= 15(40) \\
                        &= 600 \\
                  P(45) &= 600+22.5(45-40) \\
                        &= 600+22.5(5) \\
                        &= 712.50
                \end{work}
                \writelines{2}
        \end{enumerate}
  \item \textbf{Read a step (parking) graph.} A garage charges \$2 for each hour \emph{or part of an
        hour}. The cost $C(t)$ for $t$ hours parked is graphed.
        \begin{center}
        \begin{tikzpicture}[scale=0.5]
          \lgrid{0}{4}{0}{8}
          \draw[very thick, forest] (0,2)--(1,2); \opendot{0}{2} \closeddot{1}{2}
          \draw[very thick, forest] (1,4)--(2,4); \opendot{1}{4} \closeddot{2}{4}
          \draw[very thick, forest] (2,6)--(3,6); \opendot{2}{6} \closeddot{3}{6}
          \draw[very thick, forest] (3,8)--(4,8); \opendot{3}{8} \closeddot{4}{8}
        \end{tikzpicture}
        \end{center}
        $C(1.5)=\blank{1.2cm}$; \quad $C(3)=\blank{1.2cm}$. \quad Why does the graph \emph{jump} at each
        whole hour? \writelines{3}
  \item \textbf{Reason about continuity.} A piecewise function's two pieces meet with \emph{no jump} at a
        boundary $x=a$ exactly when what is true about the two rules at $a$? Explain. \writelines{3}
\end{enumerate}
\end{tcolorbox}

\end{document}
